\section{Conclusion}
\label{sec:conclusion}

We introduced the concept of \emph{natural hallucinations}---factual errors that are consistent across models, robust to rephrasing, and resistant to self-detection---and provided systematic empirical evidence for their existence.
Testing four OpenAI models on all 817 \truthfulqa questions, we found 27 universal hallucinations that fool every model, a 76\% persistence rate under rephrasing, self-detection accuracy of only 10--16\% (below chance), and cross-model Jaccard overlap 3--6$\times$ above random.
Older-model hallucinations predict newer-model hallucinations with correlations up to $r = 0.53$.

These results suggest that certain hallucinations are not bugs to be patched through incremental improvements, but structural features of how language models learn from data that contains widely-believed misconceptions.
Addressing natural hallucinations will likely require targeted interventions---explicit correction in training data, retrieval-augmented generation, or external verification systems---rather than relying on scale or self-critique alone.

Future work should extend this analysis to models from different providers to test true cross-family transfer, investigate whether retrieval augmentation can break natural hallucinations at inference time, and use mechanistic interpretability tools on open-source models to understand \emph{why} internal knowledge fails to surface for these specific questions.
